% Источники: exam/materials/прошлогодние ответы/Modelirovanie_1_-3.pdf, вопрос 8; https://github.com/HumsterProgrammer/iu9-notes/blob/main/8sem/Моделирование/Моделирование_лекция-04_19-02-26.md; https://github.com/HumsterProgrammer/iu9-notes/blob/main/8sem/Моделирование/Моделирование_лекция-05_26-02-26.md.
\section{Подход к имитации на основе процессов и транзактов.}

Под квазипараллелизмом понимают организацию моделирования объектов или систем, параллельно функционирующих в реальном мире, на компьютерах последовательной архитектуры.

\subsection*{Транзактный подход}

\textbf{Транзакт} -- заявка на обслуживание. Транзакты имеют индивидуальный закон распределения; закон распределения случайной величины (транзакта) может быть выбран из анализа предметной области. Если распределение неизвестно, необходимо провести натурный эксперимент, поскольку ошибка в выборе распределения приводит к ошибке модели.

В лекции транзакт задается как
\[
  T_i = \{D_i,\; \sum_k \tau^k_{ij}\},
\]
где $D_i$ -- дисциплина обработки заявки, а сумма времен задает время ожидания, обработки и деинициации заявки.

Генерация заявок осуществляется устройствами: станками, рабочими станциями, кассами, окнами регистратуры и т.п. Поступающие заявки организуются в очереди и регламентируются семафором. Если семафор равен 1, выполняется первая вошедшая в очередь заявка либо заявка с высшим приоритетом.

\subsection*{Процессный подход}

\textbf{Процесс} -- неделимая единица исследования, действие. В процессной организации квазипараллелизма в качестве заявки выступает не объект, а запрос на действие. Рассматривается дисциплина обработки таких запросов, существующих в данное время.

В лекции процесс задается как
\[
  P_i = \{P_i^*,\; \sum_k \tau^k_{ij}\},
\]
где $P_i^*$ -- приближенное описание процесса (иногда алгоритм), а сумма времен характеризует реальное или модельное время выполнения частей процесса.

Пример процессного подхода из лекции: получение направления -- подпись заместителя декана, подпись заведующего кафедрой, сдача задания, оценка преподавателя, учет направления и возврат в деканат. Объекты в такой модели простые, но процессы сложны и являются главным элементом модели. Обработка может быть сведена к транзактному подходу, если считать транзактом запрос на действие.

\clearpage
